Between 1950 and 1975 millions of Spaniards left the villages of the interior for Madrid, Barcelona and other cities. People in Spain still speak of going back to the village on holiday. I ask whether this shows in where people actually sleep. From mobile-phone estimates of overnight stays published by the Spanish transport ministry, I build one directed network per night of 2024: the nodes are 2,435 zones of mainland Spain, and a link goes from zone $i$ to zone $j$ when at least 10 residents of $i$ sleep in $j$. I compare an ordinary weekday in February with the summer holiday of 14--15 August and two family holidays, and set these networks against a third one built from the census: where the residents of each city were born. On the summer night the network is almost four times denser and splits into zones that send people and zones that receive them; random networks with the same degrees reproduce most of the change in how its nodes mix. On every night, a city sends more of its residents to the provinces where more of them were born. In August this pull grows, but only towards inland destinations and only through residents who arrived before 1981 and those born in villages. In the fifth of zones that lost most population, 24\,\% receive more people than they send on the ordinary night and 99\,\% on 15 August. The beaches take a large part of the summer; the inland part follows the map of the exodus.
